\section{结论与展望}

本文针对车辆轨迹优化中的噪声不确定性与遮挡干扰，提出了一种融合自适应协方差估计与野值剔除的卡尔曼滤波框架。实验表明，该方法在RMSE与轨迹平滑性上均有显著改进，且通过在线调整噪声参数，避免了人工调参带来的经验依赖。

当前方法主要依赖于线性运动模型，在高度非线性的急转弯场景下仍存在局限性。未来工作将探索基于无迹卡尔曼滤波（UKF）或粒子滤波的非线性扩展方案，并结合多源传感器融合（如IMU与视觉里程计）以进一步增强复杂工况下的轨迹恢复能力。此外，面向大规模无标签数据的自监督轨迹平滑将成为另一值得深入的方向。